\newcommand{\conone}{\\
The shell script successfully identifies and prints all prime numbers within a user-defined range using command-line arguments. It demonstrates the use of conditional statements, loops, and modular arithmetic in shell scripting. This task not only reinforces fundamental programming concepts but also enhances the ability to write efficient and reusable scripts for numerical problems in Unix/Linux environments.
}

\newcommand{\contwo}{\\
This shell script efficiently prints Pascal’s Triangle using loops and arithmetic expressions. By implementing the recurrence relation for binomial coefficients, it avoids the complexity of calculating factorials and enhances performance. The assignment demonstrates key programming concepts like nested loops, conditional statements, and mathematical logic in shell scripting, reinforcing the understanding of combinatorics and algorithmic thinking.
}

\newcommand{\conthree}{\\
This shell script checks for a valid filename argument and ensures the file exists. It then deletes all lines containing the word \texttt{college} from the file using \texttt{sed}, and confirms the deletion with a success message.
}

\newcommand{\confour}{\\
This shell script accepts a filename and two line numbers \texttt{m} and \texttt{n} as arguments. It verifies that the file exists and that the line numbers are valid positive integers with \texttt{m} less than or equal to \texttt{n}. It then prints the lines from \texttt{m} to \texttt{n} from the specified file using the \texttt{sed} command.
}

\newcommand{\confive}{\\
This shell script checks if exactly one filename is provided as an argument and verifies that the file exists. It then uses the \texttt{grep} command to display all lines from the file that start with a vowel (case-insensitive).
}

\newcommand{\consix}{\\
This shell script processes each argument passed to it, checks whether it is a regular file or a directory, and displays appropriate information. For regular files, it also counts and prints the number of lines using the \texttt{wc -l} command. If the item is neither a file nor a directory, it prints a corresponding message.
}

\newcommand{\conseven}{\\
This shell script checks if exactly one directory name is provided and verifies whether it exists. If valid, it proceeds to identify empty subfolders within the given directory and indicates that their names are listed in the file \texttt{empty\_subfolders.txt}.
}

\newcommand{\coneight}{\\
This shell script verifies that a filename is provided and the file exists. It then calculates and displays the number of lines, words, and characters in the file using the \texttt{wc} command.
}

\newcommand{\connine}{\\
This shell script reads a list of numbers from the user and sorts them using the bubble sort algorithm. It takes one command-line argument to determine the sorting order: \texttt{1} for ascending and \texttt{2} for descending. The script then displays the sorted array based on the selected order.
}

\newcommand{\conten}{\\
This shell script generates and displays the Fibonacci series up to a given range \texttt{n}, provided as a command-line argument. It uses a \texttt{while} loop to compute the series iteratively and prints each term until the value exceeds \texttt{n}.
}

\newcommand{\coneleven}{\\
This shell script takes a list of numbers as input from the user and identifies the second highest number. It ensures that at least two distinct numbers are entered and handles cases where all numbers may be the same.
}

\newcommand{\contwelve}{\\
This shell script performs a binary search on a sorted list of numbers entered by the user. It efficiently searches for a given number and reports its index if found, or a message indicating it was not found.
}

\newcommand{\conthirteen}{\\
This shell script reads a text file and generates a sorted list of unique words along with their frequencies, in descending order. It normalizes the text by converting it to lowercase, replaces punctuation and spaces with newlines, removes empty lines, and uses \texttt{sort}, \texttt{uniq -c}, and \texttt{sort -nr} to compute and display the word frequency.
}

\newcommand{\confourteen}{\\
This shell script accepts a filename as a command-line argument, verifies its existence, and removes duplicate lines by sorting the file and using the \texttt{uniq} command. The resulting unique lines are saved to \texttt{unique\_lines.txt}.
}

\newcommand{\confifteen}{\\
This shell script takes an input and an output file as arguments, checks the existence of the input file, and removes all vowels (both uppercase and lowercase) from its content using the \texttt{tr} command. The result is saved to the specified output file.
}

\newcommand{\consixteen}{\\
This shell script solves a quadratic equation of the form \( ax^2 + bx + c = 0 \). It first checks whether the coefficient \( a \) is non-zero, then calculates the discriminant to determine the nature of the roots. Based on the discriminant, it computes and displays either two real and distinct roots, one real repeated root, or two complex roots.
}

\newcommand{\conseventeen}{\\
This shell script sorts a list of user-entered numbers in descending order using the selection sort algorithm. It repeatedly selects the maximum element from the unsorted part and places it at the beginning. The sorted array is then displayed.
}

\newcommand{\coneighteen}{\\
This shell script calculates the Greatest Common Divisor (GCD) of two user-provided numbers using the Euclidean algorithm. It also handles negative inputs by converting them to their absolute values before computation.
}

\newcommand{\connineteen}{\\
This shell script performs basic arithmetic operations—addition, subtraction, multiplication, and division—based on the user's choice provided as a command-line argument. It reads two numbers from the user and uses a \texttt{case} statement to execute the selected operation, with proper handling for division by zero and invalid input.
}

\newcommand{\contwenty}{\\
This shell script checks whether a given number is a palindrome. It reverses the input number and compares it with the original. If both are equal, the number is a palindrome; otherwise, it is not.
}

\newcommand{\contwentyone}{\\
This shell script allows the user to convert the contents of a given file to either uppercase or lowercase. It verifies the file’s existence and uses the \texttt{tr} command based on the user's choice to perform the case conversion.
}

\newcommand{\contwentytwo}{\\
This shell script sorts a list of numbers entered by the user in ascending order using the bubble sort algorithm. It repeatedly compares and swaps adjacent elements to arrange the array in the correct order, and finally displays the sorted result.
}

\newcommand{\contwentythree}{\\
This shell script prints a diamond-like star pattern using nested \texttt{for} loops. It first prints an increasing sequence of stars (from 1 to 4), followed by a decreasing sequence (from 3 to 1), creating a symmetric pattern.
}

\newcommand{\contwentyfour}{\\
This shell script computes and displays the prime factors of a given positive integer greater than 1. It validates the input and uses a loop to divide the number by the smallest possible factor iteratively until all prime factors are found.
}

\newcommand{\contwentyfive}{\\
This shell script takes a filename and a word as input, checks whether the file exists, and then counts the number of times the exact word appears in the file using the \texttt{grep} and \texttt{wc} commands. It ensures that only whole word matches are counted.
}

\newcommand{\contwentysix}{\\
This shell script checks whether a user-entered word (with a minimum length of 4 characters) is a palindrome. It compares the word with its reverse and displays an appropriate message based on the result.
}

\newcommand{\contwentyseven}{\\
This shell script takes a filename as input, verifies its existence, and then calculates and displays the number of lines, words, and characters in the file using the \texttt{wc} command.
}

\newcommand{\contwentyeight}{\\
This shell script retrieves the current hour and displays an appropriate greeting based on the time of day. It categorizes the time into morning, afternoon, evening, or night using conditional checks.
}

\newcommand{\contwentynine}{\\
This shell script performs a linear search on an array of numbers entered by the user. It checks each element to find the specified key and prints its position if found, otherwise it notifies that the number is not present in the array.
}

\newcommand{\conthirty}{\\
This shell script implements the binary search algorithm on a sorted array provided by the user. It efficiently searches for a specified number and reports its position if found; otherwise, it informs the user that the number is not present in the array.
}

\newcommand{\conthirtyone}{\\
This shell script calculates the factorial of a non-negative integer entered by the user. It validates the input and uses a \texttt{for} loop to compute the factorial iteratively, then displays the result.
}

\newcommand{\conthirtytwo}{\\
This shell script checks whether a given string is a \textbf{palindrome}. It uses the \texttt{rev} command to reverse the input string and compares the reversed string with the original. If both strings match, it confirms the input is a palindrome; otherwise, it states it is not.

This program demonstrates a simple yet effective use of string manipulation and conditional statements in shell scripting, making it a helpful exercise for beginners learning Bash scripting.
}

\newcommand{\conthirtythree}{\\
This shell script implements a basic \textbf{Employee Record Management System} using a command-line interface. It allows the user to:

\begin{itemize}
    \item Insert new employee records into a file (\texttt{employee\_data.txt}).
    \item Display all existing employee records.
    \item Search and display the salary of an employee using their ID.
    \item Exit the program gracefully.
\end{itemize}

The script uses standard Bash constructs such as loops, conditionals, and file handling operations. It also employs common utilities like \texttt{grep} and \texttt{cut} to manipulate and extract data. The system ensures data persistence by storing records in a file and handles invalid inputs effectively.

Overall, this project is a simple yet effective way to understand basic shell scripting, especially for beginners interested in learning file processing and menu-driven programming in Unix/Linux environments.
}